% =====================================================================
%  Badania eksperymentalne — co sprawdzono, jak i z jakim skutkiem:
%  metodyka, zbiory danych, wyniki (tabele/wykresy) i ich interpretacja.
%  Wynik bez interpretacji to nie badanie, tylko screenshot.
%
%  Demonstracja tabeli: tabular + caption + label (label PO caption).
% =====================================================================
\section{Badania eksperymentalne i walidacja biblioteki}\label{sec:badania}
Tutaj coś piszemy...
\subsection{Zbiór XOR}
Zbiór danych XOR składa się z czterech tensorów o wymiarach
$2 \times 1$, zawierających wszystkie możliwe kombinacje wartości
dla funkcji logicznej XOR - co ukazano w tabeli
\ref{tab:xor-dataset}
(str. \pageref{tab:xor-dataset}).
\begin{table}[H]
	\centering
	\caption{Zbiór danych XOR. Źródło: opracowanie własne.}
	\begin{tabular}{|c|c|}
		\hline
		Wejście      & Wyjście \\
		\hline
		$[0, 0]^{T}$ & $[0]$   \\
		$[0, 1]^{T}$ & $[1]$   \\
		$[1, 0]^{T}$ & $[1]$   \\
		$[1, 1]^{T}$ & $[0]$   \\
		\hline
	\end{tabular}

	\label{tab:xor-dataset}
\end{table}
